

% Este trabajo de investigación es origen de varias publicaciones científicas, las dos primeras son requisito del Programa de Doctorado en Educación de la Universidad de Murcia,  al que pertenece, la tercera es adicional.


% \textbf{Congreso Internacional}

% Arteaga Marin, M. I., Olivares Carrillo, P. y Maurandi López, A. (2022). Metodologías activas en el proceso educativo: recomendaciones para una
% implementación eficaz en educación básica y bachillerato para la enseñanza
% STEM. \textit{3rd edition of Porto International Conference on Research in Education},
% \url{https://porto-icre2022.eventqualia.net/en/2022/home/}.

% \textbf{Artículos Indexados}

% \textcolor{red}{Latindex}. Arteaga-Marín, M. I. ., Sánchez-Rodríguez, A. ., Olivares-Carrillo,P., y Maurandi-López, A. (2022). Revisión sistemática y propuesta para la implementación de metodologías activas en la educación STEM. \textit{EDUCATECONCIENCIA, 30}(36), 35-76, \url{https://tecnocientifica.com.mx/educateconciencia/index.php/revistaeducate/article/view/533}. ISSN: 2683-2836

% \textcolor{red}{Scopus Q1}. Ibáñez-López, F. J.; Arteaga-Marín, M.; Olivares-Carrillo, P.; Sánchez-Rodríguez, A.; Maurandi-López, A. (2022). Diseño y validación de un cuestionario sobre uso de herramientas tecnológicas en innovación de asignaturas STEM. \textit{Campus Virtuales, 11}(2), 179-195, \url{https://doi.org/10.54988/cv.2022.2.1081. ISSN: 2255-1514}

% \begin{center}
%     \includegraphics[width=0.8\textwidth]{pics/publicaciones.png}
% \end{center}

%En este trabajo emplearemos las siguientes convenciones: 

%\begin{itemize}
%\item 
%Se empleará la tipografía \textit{cursiva} para palabras en otro idioma distinto del castellano y para resaltar nombres como, por ejemplo, los nombres de variables.
%\item 
%Emplearemos la tipografía \texttt{código} para sentencias de código,  nombres de funciones y librerías de \texttt{R}.
%\item 
%Emplearemos el punto como separador decimal, por ejemplo: \texttt{1/2=0.5}, y un espacio simple como separador de millares, por ejemplo: \texttt{123\;345} tal como recomienda la Academia de la Lengua Española en la Ortografía del 2010, página 666 \cite{espanola2010ortografia}. 
%\item 
%Cuando informemos sobre una media, la acompañaremos entre paréntesis de la desviación típica, por ejemplo: \texttt{8.81(1.21)}. 
%\item 
%Si informamos de un porcentaje, entre paréntesis, incluimos la frecuencia absoluta, por ejemplo: \texttt{50.25\%(123)}, o viceversa: \texttt{123(50.25\%)}. 
%\item 
%En las tablas de descriptivos, empleamos la abreviatura \texttt{SD} para \emph{desviación estandar} (del inglés \emph{Standard  Deviation}).
%\item 
%En el texto empleamos la abreviatura \texttt{EPV} para \emph{Enfermedades Prevenibles por Vacunas}.
%\end{itemize}

